% !PRINTED 635
\weberpage{635}
%<<CONT
\weq{5}{\begin{aligned}
(\varepsilon^{g}, x)(\varepsilon, x)^{-g} &= f_0\\
(\varepsilon^{g^2}, x)(\varepsilon^{g}, x)^{-g} &= f_1\\
\cdot\ \cdot\ \cdot\ \cdot\ \cdot\ \cdot\ \cdot\ \cdot\ \cdot\ \cdot & \\
(\varepsilon^{g^{n-1}}, x)(\varepsilon^{g^{n-2}}, x)^{-g} &= f_{n-2}.
\end{aligned}}
so that, if $(\varepsilon^{g}, x)(\varepsilon, x)^{-g} = f(\varepsilon)$ is set, we have
\weq{6}{f_h = f(\varepsilon^{g^h})}
If one applies the permutation $t^{-1}$ to (5), then, as is to be seen from Theorem 4., the functions
\[
f_0, f_1, f_2 \ldots f_{n-2}
\]
undergo a cyclic permutation $(0, 1, 2 \ldots n-2)$.

Moreover, so long as the $x$ are variable, the functions $f_\nu$ are, as one sees from (5), all \gesperrt{different from one another} since they are decomposed into different linear factors.

If, then, we form a cyclic function of these $n-1$ quantities
\weq{7}{C(f_0, f_1, f_2 \ldots f_{n-2}),}
this remains unchanged both by $s$ and by $t^{-1}$, hence by $t$ and the whole metacyclic group, and we obtain the theorem:

\gesperrt{6. If a cyclic function of the quantities $f_0, f_1 \ldots f_{n-1}$ is represented in the normal form, then the coefficients are metacyclic functions of the variables $x_0, x_1 \ldots x_{n-1}$.} % [sic] Print says $f_{n-1}$ although (5) defines the sequence only through $f_{n-2}$; preserved.

Here a metacyclic function is to be understood as any function that remains unchanged by the permutations of the linear group.

But the functions $f_0, f_1 \ldots f_{n-2}$ also pass cyclically into one another if one leaves the $x$ unchanged and makes the substitution $(\varepsilon, \varepsilon^g)$ for $\varepsilon$. If, then, the cyclic function $C$ of the $f$ does not contain $\varepsilon$ in its coefficients, then $C$ remains unchanged by all the substitutions $(\varepsilon, \varepsilon^h)$, and we obtain from 3. the theorem:

\gesperrt{7. If $C(f_0, f_1 \ldots f_{n-2})$ is a cyclic function of the quantities $f_0, f_1 \ldots f_{n-2}$ with rational coefficients, then $C$ is a metacyclic function of the variables $x_0, x_1 \ldots x_{n-1}$ with rational coefficients.}

If we raise the functions $f_0, f_1 \ldots f_{n-2}$, as they are given by (5), in order to the powers $g^{n-2}, g^{n-3} \ldots 1$
%CONT>>
